\documentclass[journal]{IEEEtran}
\usepackage[UTF8]{ctex}
\usepackage{amsmath,amssymb,amsfonts}
\usepackage{graphicx}
\usepackage{cite}
\usepackage{algorithmic}
\usepackage{algorithm}
\usepackage{array}
\usepackage{url}
\usepackage{booktabs}
\usepackage{multirow}
\usepackage{bm}

\hyphenation{op-tical net-works semi-conduc-tor}

\begin{document}

\title{Anytime-Valid Sequential Testing for Robust UAV Gas Source Localization with Distribution-Free Stopping Guarantees}

\author{
\IEEEauthorblockN{Author Name}
\IEEEauthorblockA{Affiliation}
}

\markboth{IEEE Sensors Journal}%
{Author: Anytime-Valid UAV Gas Source Localization}

\maketitle

\begin{abstract}
Unmanned aerial vehicles (UAVs) for gas source localization (GSL) face a critical challenge: classical upwind-tracking methods depend on continuous plume contact and fail when the initial gas signal is weak, with localization error increasing up to 8$\times$. Moreover, no existing GSL method provides formal guarantees on when to stop searching. We propose a unified framework based on Blackwell's theory of approachability that jointly optimizes search efficiency, stopping decisions, and risk sensitivity. For the stopping criterion, we construct an anytime-valid e-process from posterior convergence indicators, providing distribution-free false termination guarantees ($\mathbb{P}(\text{false positive}) \leq \delta$) without requiring exchangeability assumptions. For search strategy, we employ risk-adjusted information-directed sampling that balances exploration and exploitation. We prove that: (i) the target coverage-efficiency set is approachable under mild identifiability conditions; (ii) the e-process stopping rule satisfies anytime-validity via Ville's inequality; and (iii) the risk-adjusted search accelerates posterior convergence. Experiments on three indoor environments with varying signal strengths demonstrate consistent localization performance (1.2--2.1m mean distance) across all environments, with 2$\times$ lower cross-environment variance than the state-of-the-art SurgeCast method. The e-process stopping criterion fires in 100\% of runs (vs.\ 0\% for previous conformal prediction approaches), providing actionable confidence for deployment.
\end{abstract}

\begin{IEEEkeywords}
Gas source localization, unmanned aerial vehicles, anytime-valid inference, e-processes, Blackwell approachability, sequential hypothesis testing
\end{IEEEkeywords}

\input{sections/introduction}
\input{sections/related_work}
\input{sections/methodology}
\input{sections/experiments}
\input{sections/references}

\end{document}